%!TeX root=../gardentop.tex
\chapter{The Nest of the Missel Thrush} 
	
\begin{figure}[t!]
\centering
\includegraphics[width=\linewidth]{january}
\end{figure}

 \lettrine[lines=6]{F}{or} two or three minutes he stood looking round him, while Mary watched him, and then he began to walk about softly, even more lightly than Mary had walked the first time she had found herself inside the four walls. His eyes seemed to be taking in everything—the gray trees with the gray creepers climbing over them and hanging from their branches, the tangle on the walls and among the grass, the evergreen alcoves with the stone seats and tall flower urns standing in them.

»I never thought I'd see this place,« he said at last, in a whisper.

»Did you know about it?« asked Mary.

She had spoken aloud and he made a sign to her.

»We must talk low,« he said, »or some one'll hear us an' wonder what's to do in here.«

»Oh! I forgot!« said Mary, feeling frightened and putting her hand quickly against her mouth. »Did you know about the garden?« she asked again when she had recovered herself.

Dickon nodded.

»Martha told me there was one as no one ever went inside,« he answered. »Us used to wonder what it was like.«

He stopped and looked round at the lovely gray tangle about him, and his round eyes looked queerly happy.

»Eh! the nests as'll be here come springtime,« he said. »It'd be th' safest nestin' place in England. No one never comin' near an' tangles o' trees an' roses to build in. I wonder all th' birds on th' moor don't build here.«

Mistress Mary put her hand on his arm again without knowing it.

»Will there be roses?« she whispered. »Can you tell? I thought perhaps they were all dead.«

»Eh! No! Not them—not all of 'em!« he answered. »Look here!«

He stepped over to the nearest tree—an old, old one with gray lichen all over its bark, but upholding a curtain of tangled sprays and branches. He took a thick knife out of his pocket and opened one of its blades.

»There's lots o' dead wood as ought to be cut out,« he said. »An' there's a lot o' old wood, but it made some new last year. This here's a new bit,« and he touched a shoot which looked brownish green instead of hard, dry gray.

Mary touched it herself in an eager, reverent way.

»That one?« she said. »Is that one quite alive—quite?«

Dickon curved his wide smiling mouth.

»It's as wick as you or me,« he said; and Mary remembered that Martha had told her that »wick« meant »alive« or »lively.«

»I'm glad it's wick!« she cried out in her whisper. »I want them all to be wick. Let us go round the garden and count how many wick ones there are.«

She quite panted with eagerness, and Dickon was as eager as she was. They went from tree to tree and from bush to bush. Dickon carried his knife in his hand and showed her things which she thought wonderful.

»They've run wild,« he said, »but th' strongest ones has fair thrived on it. The delicatest ones has died out, but th' others has growed an' growed, an' spread an' spread, till they's a wonder. See here!« and he pulled down a thick gray, dry-looking branch. »A body might think this was dead wood, but I don't believe it is—down to th' root. I'll cut it low down an' see.«

He knelt and with his knife cut the lifeless-looking branch through, not far above the earth.

»There!« he said exultantly. »I told thee so. There's green in that wood yet. Look at it.«

Mary was down on her knees before he spoke, gazing with all her might.

»When it looks a bit greenish an' juicy like that, it's wick,« he explained. »When th' inside is dry an' breaks easy, like this here piece I've cut off, it's done for. There's a big root here as all this live wood sprung out of, an' if th' old wood's cut off an' it's dug round, an' took care of there'll be—« he stopped and lifted his face to look up at the climbing and hanging sprays above him—»there'll be a fountain o' roses here this summer.«

They went from bush to bush and from tree to tree. He was very strong and clever with his knife and knew how to cut the dry and dead wood away, and could tell when an unpromising bough or twig had still green life in it. In the course of half an hour Mary thought she could tell too, and when he cut through a lifeless-looking branch she would cry out joyfully under her breath when she caught sight of the least shade of moist green. The spade, and hoe, and fork were very useful. He showed her how to use the fork while he dug about roots with the spade and stirred the earth and let the air in.

They were working industriously round one of the biggest standard roses when he caught sight of something which made him utter an exclamation of surprise.

»Why!« he cried, pointing to the grass a few feet away. »Who did that there?«

It was one of Mary's own little clearings round the pale green points.

»I did it,« said Mary.

»Why, I thought tha' didn't know nothin' about gardenin',« he exclaimed.

»I don't,« she answered, »but they were so little, and the grass was so thick and strong, and they looked as if they had no room to breathe. So I made a place for them. I don't even know what they are.«

Dickon went and knelt down by them, smiling his wide smile.

»Tha' was right,« he said. »A gardener couldn't have told thee better. They'll grow now like Jack's bean-stalk. They're crocuses an' snowdrops, an' these here is narcissuses,« turning to another patch, »an' here's daffydowndillys. Eh! they will be a sight.«

He ran from one clearing to another.

»Tha' has done a lot o' work for such a little wench,« he said, looking her over.

»I'm growing fatter,« said Mary, »and I'm growing stronger. I used always to be tired. When I dig I'm not tired at all. I like to smell the earth when it's turned up.«

»It's rare good for thee,« he said, nodding his head wisely. »There's naught as nice as th' smell o' good clean earth, except th' smell o' fresh growin' things when th' rain falls on 'em. I get out on th' moor many a day when it's rainin' an' I lie under a bush an' listen to th' soft swish o' drops on th' heather an' I just sniff an' sniff. My nose end fair quivers like a rabbit's, mother says.«

»Do you never catch cold?« inquired Mary, gazing at him wonderingly. She had never seen such a funny boy, or such a nice one.

»Not me,« he said, grinning. »I never ketched cold since I was born. I wasn't brought up nesh enough. I've chased about th' moor in all weathers same as th' rabbits does. Mother says I've sniffed up too much fresh air for twelve year' to ever get to sniffin' with cold. I'm as tough as a white-thorn knobstick.«

He was working all the time he was talking and Mary was following him and helping him with her fork or the trowel.

»There's a lot of work to do here!« he said once, looking about quite exultantly.

»Will you come again and help me to do it?« Mary begged. »I'm sure I can help, too. I can dig and pull up weeds, and do whatever you tell me. Oh! do come, Dickon!«

»I'll come every day if tha' wants me, rain or shine,« he answered stoutly. »It's th' best fun I ever had in my life—shut in here an' wakenin' up a garden.«

»If you will come,« said Mary, »if you will help me to make it alive I'll—I don't know what I'll do,« she ended helplessly. What could you do for a boy like that?

»I'll tell thee what tha'll do,« said Dickon, with his happy grin. »Tha'll get fat an' tha'll get as hungry as a young fox an' tha'll learn how to talk to th' robin same as I do. Eh! we'll have a lot o' fun.«

He began to walk about, looking up in the trees and at the walls and bushes with a thoughtful expression.

»I wouldn't want to make it look like a gardener's garden, all clipped an' spick an' span, would you?« he said. »It's nicer like this with things runnin' wild, an' swingin' an' catchin' hold of each other.«

»Don't let us make it tidy,« said Mary anxiously. »It wouldn't seem like a secret garden if it was tidy.«

Dickon stood rubbing his rusty-red head with a rather puzzled look.

»It's a secret garden sure enough,« he said, »but seems like some one besides th' robin must have been in it since it was shut up ten year' ago.«

»But the door was locked and the key was buried,« said Mary. »No one could get in.«

»That's true,« he answered. »It's a queer place. Seems to me as if there'd been a bit o' prunin' done here an' there, later than ten year' ago.«

»But how could it have been done?« said Mary.

He was examining a branch of a standard rose and he shook his head.

»Aye! how could it!« he murmured. »With th' door locked an' th' key buried.«

Mistress Mary always felt that however many years she lived she should never forget that first morning when her garden began to grow. Of course, it did seem to begin to grow for her that morning. When Dickon began to clear places to plant seeds, she remembered what Basil had sung at her when he wanted to tease her.

»Are there any flowers that look like bells?« she inquired.

»Lilies o' th' valley does,« he answered, digging away with the trowel, »an' there's Canterbury bells, an' campanulas.«

»Let us plant some,« said Mary.

»There's lilies o' th' valley here already; I saw 'em. They'll have growed too close an' we'll have to separate 'em, but there's plenty. Th' other ones takes two years to bloom from seed, but I can bring you some bits o' plants from our cottage garden. Why does tha' want 'em?«

Then Mary told him about Basil and his brothers and sisters in India and of how she had hated them and of their calling her »Mistress Mary Quite Contrary.«

»They used to dance round and sing at me. They sang—

\begin{verse}
Mistress Mary, quite contrary,\\
How does your garden grow?\\
With silver bells, and cockle shells,\\
And marigolds all in a row.
\end{verse}

\noindent I just remembered it and it made me wonder if there were really flowers like silver bells.«

She frowned a little and gave her trowel a rather spiteful dig into the earth.

»I wasn't as contrary as they were.«

But Dickon laughed.

»Eh!« he said, and as he crumbled the rich black soil she saw he was sniffing up the scent of it, »there doesn't seem to be no need for no one to be contrary when there's flowers an' such like, an' such lots o' friendly wild things runnin' about makin' homes for themselves, or buildin' nests an' singin' an' whistlin', does there?«

Mary, kneeling by him holding the seeds, looked at him and stopped frowning.

»Dickon,« she said. »You are as nice as Martha said you were. I like you, and you make the fifth person. I never thought I should like five people.«

Dickon sat up on his heels as Martha did when she was polishing the grate. He did look funny and delightful, Mary thought, with his round blue eyes and red cheeks and happy looking turned-up nose.

»Only five folk as tha' likes?« he said. »Who is th' other four?«

»Your mother and Martha,« Mary checked them off on her fingers, »and the robin and Ben Weatherstaff.«

Dickon laughed so that he was obliged to stifle the sound by putting his arm over his mouth.

»I know tha' thinks I'm a queer lad,« he said, »but I think tha' art th' queerest little lass I ever saw.«

Then Mary did a strange thing. She leaned forward and asked him a question she had never dreamed of asking any one before. And she tried to ask it in Yorkshire because that was his language, and in India a native was always pleased if you knew his speech.

»Does tha' like me?« she said.

»Eh!« he answered heartily, »that I does. I likes thee wonderful, an' so does th' robin, I do believe!«

»That's two, then,« said Mary. »That's two for me.«

And then they began to work harder than ever and more joyfully. Mary was startled and sorry when she heard the big clock in the courtyard strike the hour of her midday dinner.

»I shall have to go,« she said mournfully. »And you will have to go too, won't you?«

Dickon grinned.

»My dinner's easy to carry about with me,« he said. »Mother always lets me put a bit o' somethin' in my pocket.«

He picked up his coat from the grass and brought out of a pocket a lumpy little bundle tied up in a quiet clean, coarse, blue and white handkerchief. It held two thick pieces of bread with a slice of something laid between them.

»It's oftenest naught but bread,« he said, »but I've got a fine slice o' fat bacon with it to-day.«

Mary thought it looked a queer dinner, but he seemed ready to enjoy it.

»Run on an' get thy victuals,« he said. »I'll be done with mine first. I'll get some more work done before I start back home.«

He sat down with his back against a tree.

»I'll call th' robin up,« he said, »and give him th' rind o' th' bacon to peck at. They likes a bit o' fat wonderful.«

Mary could scarcely bear to leave him. Suddenly it seemed as if he might be a sort of wood fairy who might be gone when she came into the garden again. He seemed too good to be true. She went slowly half-way to the door in the wall and then she stopped and went back.

»Whatever happens, you—you never would tell?« she said.

His poppy-coloured cheeks were distended with his first big bite of bread and bacon, but he managed to smile encouragingly.

»If tha' was a missel thrush an' showed me where thy nest was, does tha' think I'd tell any one? Not me,« he said. »Tha' art as safe as a missel thrush.«

And she was quite sure she was.